

\section{Methodology}



%%%%%%%%%%%%%%%%%%%%%%%%%%%%%%%%%%%%%%%%%%%
\subsection{Continual Learning}
Given a set of $T$ tasks $\{{\mathcal{T}^t}\}_{t=1}^{T}$, continual learning works by sequentially accessing and learning on each task $\mathcal{T}^t=\{\mathcal{D}^t,\mathcal{C}^t\}$. Here, $\mathcal{D}^t=\{{I}^{t}_i,{y}^{t}_i\}_{i=1}^{N^t}$ represents the data of $t^{th}$ task $\mathcal{T}^t$, where $I^{t}_{i}$ is the input image, $y^{t}_{i} \in \mathcal{C}^t$ is the corresponding class label, and $N^t$ is the size of data. The category set $\mathcal{C}^t=\{c^{t}_j\}_{j=1}^{M^t}$ encompasses the class names within $\mathcal{T}^t$, with a total of $M^t$ classes.
Continual learning aims to achieve good performance across all tasks and can be broadly categorized into Task Incremental Learning (TIL) and Class Incremental Learning (CIL).  In TIL, the model generates predictions within a task-specific set $\mathcal{C}^t$, which is determined by the current task identity $t$. Meanwhile, in CIL, the challenge involves distinguishing between all the previously encountered classes $\cup_{i=1}^t \mathcal{C}^i$.



%%%%%%%%%%%%%%%%%%%%%%%%%%%%%%%%%%%%%%%%%%%%%%%%

\subsection{Framework Overview}


%%%%%%%%%%%%%%%%%%%%%%%%%%%%%%%%%%%%%%%%%%%%%%%%
In this paper, we present a parameter-efficient framework designed to empower the continual learning capabilities of vision-language models, achieving robust historical knowledge memorization without sacrificing the  zero-shot generalization abilities. Our method is built upon the CLIP~\cite{gao2023clip} model, which contains parallel encoders $(\mathcal{F}_I,\mathcal{F}_T)$ to extract features of input images and texts, respectively. By following  CLIP~\cite{gao2023clip}, we make predictions based on the cosine similarity between the final image embedding $\mathcal{F}_I(I_i^t)$ and text embedding $\mathcal{F}_T(c^{t}_{j})$. The input image is then assigned to the class with the highest similarity.





%%%%%%%%%%%%%%%%%%%%%%%%%%%%%%%%%%%%%%%%%%%%%%%%
The overall framework of our method is shown in Figure~\ref{fig:overview}. We introduce MoE structure onto a frozen CLIP to consolidate all the downstream tasks within a unified model, in which the task-dependent routers are sequentially added to modulate the experts for each task. Adapter modules, such as LoRA~\cite{hu2021lora}, function as the experts in the MoE setup, enhancing adaptation speed during training. To relieve MoE's reliance  on task identities, we further propose a Distribution Discriminative Auto-Selector (DDAS). The DDAS automatically infers the task context by analyzing the variations of the target image distributions.  As a result, in-distribution data will be allocated to the corresponding routers within MoE, while the out-of-distribution inputs will be identified and directed  to the original CLIP to perform zero-shot recognition.    

%-------------------------------------------------------------------------

\subsection{Incremental Mixture-of-Experts Adapters}
\label{MoEs}
%

We leverage MoE~\cite{shazeer2017outrageously} to build an expansible architecture to alleviate the ``catastrophic forgetting'' issue in the continual learning of CLIP. The MoE is composed of several experts $\{\mathcal{E}_i\}_{i=1}^{N_E}$ and routers, where $N_E$ is the number of predefined experts. For current task $\mathcal{T}^t$, only a task-dependent router $\mathcal{R}^t, t\in[1,T]$ is added to the system, which integrates the experts' outputs via gated average. 



\vspace{1ex}
\noindent\textbf{Adapters as Experts.} 
MoE in vision-language models usually incorporate the experts inside networks, which can be MLPs or attention heads~\cite{shazeer2017outrageously, zhang2022mixture, chen2023mod}. However, inserting the MoE inside VLM might bring in significant computational burdens due to the full-parameter tuning. Some methods~\cite{gao2023clip, liu2023class} have demonstrated that adapters with few parameters can increase the adaption speed of VLM on downstream tasks and enable their applications in continual learning. 
%
Inspired by this, we use the effective adapter LoRA~\cite{hu2021lora}, which works by decoupling the original heavy and frozen parameters into low-rank trainable space, as the experts in MoE to speed up continual learning with CLIP. 
%

%%%%%%%%%%%%%%%%%%%%%%%%%%%%%%%%%%%%%%%%%%%%%%%%%%%
Our MoE-Adapters are implemented in all the Transformer blocks of the parallel encoders $(\mathcal{F}_I,\mathcal{F}_T)$, as shown in Figure~\ref{fig:overview}. 
To be more specific, in each Transformer block, the feature tokens \textbf{\textit{$x^t$}} $\in \mathbb{R}^{n\times d}$ after the multi-head attention output are passed to all the experts in the MoE. 
Then, the task-specific router $\mathcal{R}^t$ is applied to fuse the experts' outputs via gated summation. Note that we implement the same MoE adapters in both the image  encoder and the textual encoder, without sharing the parameters. 
 
%
%%%%%%%%%%%%%%%%%%%%%%%%%%%%%%%%%%%%%%%%%%%%%%%%%%%



\vspace{1ex}
\noindent\textbf{Incremental Mixture of Experts.} 
%
In our MoE framework, task-specific routers $\mathcal{R}^t$ determine the activation of experts $\mathcal{E}_i$ to produce outcomes tailored to each task $t$. The combined output for a task, $\textbf{\textit{y}}^t$, is computed as:
\begin{equation}
    \textbf{\textit{y}}^t = \sum_{i=1}^{N_E}{W_i^t}\mathcal{E}_i(\textbf{\textit{x}}^t),
\end{equation}
where $W^t=\{W_i^t\}_{i=1}^{N_E}$ represents the gating weights assigned by $\mathcal{R}^t$, dictating each expert's contribution. $\textbf{\textit{x}}^t$ denotes the tokens processed for task $t$, and $\textbf{\textit{y}}^t$ is the corresponding output from the MoE-Adapters, matching the shape of $\textbf{\textit{x}}^t$.
%
We refine the MoE-Adapters for continual learning with two key modifications. Unlike previous methods~\cite{riquelme2021scaling, daxberger2023mobile} that input patch or image tokens into the router, we utilize the initial token, known as the \texttt{[CLS]} token ($\textbf{\textit{c}}^t \in \mathbb{R}^{1\times d}$), to enhance processing efficiency. The gating weights are then computed as follows:
\begin{equation}
\label{eq:gate1}
W^t = Softmax(Topk(\mathcal{R}^t(\textbf{\textit{c}}^t))),
\end{equation}
where $\mathcal{R}^t$ projects $\textbf{\textit{c}}^t$ to a 1-D vector indicating each expert's likelihood of activation. The $Topk(\cdot)$ function selects the $k$ most relevant experts, while setting the rest to be $-\infty$. The $Softmax(\cdot)$ function normalizes these weights to emphasize the selected experts' contribution.

%

\begin{figure}[t]
	\centering
	\includegraphics[width=1\linewidth]{figs/figure_differentFrozen.pdf}
	\vspace{-14pt}
	\caption{The three distinct combinations among activated experts (a) both trained, (b) trainable and frozen, (c) both experts are frozen, and only the router is trainable.}
	\label{fig:different frozen}
  \vspace{-8pt}
 \end{figure}

\begin{table*}
    \centering
    \resizebox{0.95\linewidth}{!}{
    \begin{tabular}{c>{\raggedright\arraybackslash}p{3cm} >{\centering\arraybackslash}p{1cm} >{\centering\arraybackslash}p{1cm}>{\centering\arraybackslash}p{1cm} >{\centering\arraybackslash}p{1cm} >{\centering\arraybackslash}p{1cm}>{\centering\arraybackslash}p{1cm} >{\centering\arraybackslash}p{1cm} >{\centering\arraybackslash}p{1cm}>{\centering\arraybackslash}p{1cm} >{\centering\arraybackslash}p{1cm} >{\centering\arraybackslash}p{1cm} >{\centering\arraybackslash}p{1.5cm}}
        \toprule
           & {\hspace{1em}} Method & \makecell[c]{\rotatebox{90}{Aircraft~\cite{maji2013fine}}} & \makecell[c]{\rotatebox{90}{Caltech101~\cite{fei2004learning}}} & \makecell[c]{\rotatebox{90}{CIFAR100~\cite{krizhevsky2009learning}}} & \makecell[c]{\rotatebox{90}{DTD~\cite{cimpoi2014describing}}} & \makecell[c]{\rotatebox{90}{EuroSAT~\cite{helber2019eurosat}}} & \makecell[c]{\rotatebox{90}{Flowers~\cite{nilsback2008automated}}} & \makecell[c]{\rotatebox{90}{Food~\cite{bossard2014food}}} & \makecell[c]{\rotatebox{90}{MNIST~\cite{deng2012mnist}}} & \makecell[c]{\rotatebox{90}{OxfordPet~\cite{parkhi2012cats}}} & \makecell[c]{\rotatebox{90}{Cars~\cite{krause20133d}}} & \makecell[c]{\rotatebox{90}{SUN397~\cite{xiao2010sun}}} & \makecell[c]{{\textit{Average}}} \\
  
        \midrule
        
            \multirow{3}{*}{\rotatebox{90}{CLIP}}& {\hspace{1em}}Zero-shot & 24.3 & 88.4 & 68.2 & 44.6 & 54.9 & 71.0 & 88.5 & 59.4 & 89.0 & 64.7 & 65.2 & 65.3  \\
            & {\hspace{1em}}Full Fine-tune & 62.0 & 95.1 & 89.6 & 79.5 & 98.9 & 97.5 & 92.7 & 99.6 & 94.7 & 89.6 & 81.8 & 89.2  \\
            & {\hspace{1em}}Fine-tune Adapter & 56.8 & 92.6 & 89.4 & 79.0 & 98.4 & 97.0 & 92.9 & 99.2 & 94.1 & 89.1 & 82.7 & 88.3 \\
            
            \midrule\midrule
            \multirow{8}{*}{\rotatebox{90}{\textbf{Transfer}}} &{\hspace{1em}}Continual-FT & & 67.1 & 46.0 & 32.1 & 35.6 & 35.0 & 57.7 & 44.1 & 60.8 & 20.5 & 46.6 & 44.6 \\
            & {\hspace{1em}}LwF \cite{li2017learning} &  & 74.5 & 56.9 & 39.1 & \textbf{51.1} & 52.6 & 72.8 & \underline{60.6} & 75.1 & 30.3 & 55.9 & 58.9 \\
            & {\hspace{1em}}iCaRL \cite{rebuffi2017icarl} &  & 56.6 & 44.6 & 32.7 & 39.3 & 46.6 & 68.0 & 46.0 & 77.4 & 31.9 & 60.5 & 50.4 \\
            & {\hspace{1em}}LwF-VR \cite{ding2022don} & & 77.1 & 61.0 & 40.5 & 45.3 & 54.4 & 74.6 & 47.9 & 76.7 & 36.3 & 58.6 & 57.2  \\
            & {\hspace{1em}}WiSE-FT \cite{wortsman2022robust} & & 73.5 & 55.6 & 35.6 & 41.5 & 47.0 & 68.3 & 53.9 & 69.3 & 26.8 & 51.9 & 52.3  \\
            & {\hspace{1em}}ZSCL \cite{zheng2023preventing} & & 86.0 & 67.4 & \textbf{45.4} & \underline{50.4} & \underline{69.1} & 87.6 & \textbf{61.8} & 86.8 & 60.1 & \textbf{66.8} & \underline{68.1} \\
            \rowcolor{Thistle!20}
            & {\hspace{1em}}Ours$\dagger$ & & \textbf{87.9} & \textbf{68.2} & 42.2 & 41.4 & 68.7 & \textbf{88.7} & 59.4 & \textbf{89.1} & \textbf{64.5} & 64.0 & 67.4\textbf{\textcolor{NavyBlue}{ (-0.7)}}  \\
            \rowcolor{Thistle!20}
            & {\hspace{1em}}Ours &&\textbf{87.9} & \textbf{68.2} & \underline{44.4} &49.9 &\textbf{70.7} & \textbf{88.7} &59.7 &\textbf{89.1} &\textbf{64.5} &\underline{65.5} &\textbf{68.9\textcolor{Maroon}{ (+0.8)}}\\
            % ----------------------------------------------------
            \midrule 
            \multirow{8}{*}{\rotatebox{90}{\textbf{Average}}} &{\hspace{1em}}Continual-FT    & 25.5 & 81.5 & 59.1 & 53.2 & 64.7 & 51.8 & 63.2 & 64.3 & 69.7 & 31.8 & 49.7 & 55.9 \\
            & {\hspace{1em}}LwF \cite{li2017learning} & 36.3 & 86.9 & 72.0 & 59.0 & 73.7 & 60.0 & 73.6 & \underline{74.8} & 80.0 & 37.3 & 58.1 & 64.7 \\
            & {\hspace{1em}}iCaRL \cite{rebuffi2017icarl} & 35.5 & 89.2 & 72.2 & 60.6 & 68.8 & 70.0 & 78.2 & 62.3 & 81.8 & 41.2 & 62.5 & 65.7 \\
            & {\hspace{1em}}LwF-VR \cite{ding2022don} & 29.6 & 87.7 & 74.4 & 59.5 & 72.4 & 63.6 & 77.0 & 66.7 & 81.2 & 43.7 & 60.7 & 65.1  \\
            & {\hspace{1em}}WiSE-FT \cite{wortsman2022robust} & 26.7 & 86.5 & 64.3 & 57.1 & 65.7 & 58.7 & 71.1 & 70.5 & 75.8 & 36.9 & 54.6 & 60.7  \\
            & {\hspace{1em}}ZSCL \cite{zheng2023preventing} & 45.1 & \textbf{92.0} & 80.1 & 64.3 & \textbf{79.5} & 81.6 & \textbf{89.6} & \textbf{75.2} & 88.9 & 64.7 & \textbf{68.0} &75.4 \\
            \rowcolor{Thistle!20}
            & {\hspace{1em}}Ours$\dagger$& \textbf{54.3} & 91.1 & \textbf{85.1} & \textbf{69.7} & 77.5 & \textbf{84.5} & \underline{89.1} &  73.8 & \underline{89.2} & \textbf{69.0} & 65.8 & \textbf{77.2\textcolor{Maroon}{\textbf{(+1.8)}}} \\
            \rowcolor{Thistle!20}
            & {\hspace{1em}}Ours &\underline{50.2}&\underline{91.9}&\underline{83.1}&\underline{69.4}&\underline{78.9}&\underline{84.0}&\underline{89.1}&73.7&\textbf{89.3}&\underline{67.7}&\underline{66.9}&\underline{76.7}\textcolor{Maroon}{\textbf{(+1.3)}}\\
            \midrule
            % ----------------------------------------------------
            \multirow{8}{*}{\rotatebox{90}{\textbf{Last}}} &{\hspace{1em}}Continual-FT & 31.0 & 89.3 & 65.8 & 67.3 & 88.9 & 71.1 & 85.6 & \textbf{99.6} & 92.9 & 77.3 & 81.1 & 77.3 \\
            & {\hspace{1em}}LwF \cite{li2017learning} & 26.3 & 87.5 & 71.9 & 66.6 & 79.9 & 66.9 & 83.8 & \textbf{99.6} & 92.1 & 66.1 & 80.4 & 74.6 \\
            & {\hspace{1em}}iCaRL \cite{rebuffi2017icarl} & 35.8 & \textbf{93.0} & 77.0 & 70.2 & 83.3 & 88.5 & \underline{90.4} & 86.7 & \underline{93.2} & 81.2 & \underline{81.9} & 80.1 \\
            & {\hspace{1em}}LwF-VR \cite{ding2022don} & 20.5 & 89.8 & 72.3 & 67.6 & 85.5 & 73.8 & 85.7 & \textbf{99.6} & 93.1 & 73.3 & 80.9 & 76.6  \\
            & {\hspace{1em}}WiSE-FT \cite{wortsman2022robust} & 27.2 & 90.8 & 68.0 & 68.9 & 86.9 & 74.0 & 87.6 & \textbf{99.6} & 92.6 & 77.8 & 81.3 & 77.7  \\
            & {\hspace{1em}}ZSCL \cite{zheng2023preventing} & 40.6 & \underline{92.2} & 81.3 & 70.5 & 94.8 & 90.5 & \textbf{91.9} & 98.7 & \textbf{93.9} & \underline{85.3} & 80.2 & 83.6 \\
            \rowcolor{Thistle!20}
            & {\hspace{1em}}Ours$\dagger$ & \textbf{54.3} & 90.8 & \textbf{88.8} & \textbf{80.3} & \textbf{98.1} & \textbf{97.5} & 89.6 & 99.1 & 89.5 & \textbf{89.2} & \textbf{83.8} & \textbf{87.4\textcolor{Maroon}{ (+3.8)}} \\ 
            \rowcolor{Thistle!20}
            & {\hspace{1em}}Ours& \underline{49.8}&\underline{92.2}&\underline{86.1}&\underline{78.1}&\underline{95.7}&\underline{94.3}&89.5&98.1&89.9&81.6&80.0&\underline{85.0}\textcolor{Maroon}{\textbf{(+1.4)}}\\
        \bottomrule
    \end{tabular}}
    \vspace{-4pt}
    \caption{Comparison with state-of-the-art methods on MTIL benchmark in terms of ``Transfer'', ``Average'', and ``Last'' scores (\%). ``Ours$\dagger$'' and ``Ours'' indicate our method trained on 3k and 1k iterations, respectively. We label the best and second methods with \textbf{bold} and \underline{underline} styles. The top block indicates the upper-bound solutions to adapt the CLIP on each task.}
    \label{tab:compareAvg_Last_Transfer}
     \vspace{-11pt}
\end{table*}

\begin{table*}[t]
	\centering
        \resizebox{0.99\linewidth}{!}{
	\begin{tabular}{cl>{\centering\arraybackslash}p{1cm} >{\centering\arraybackslash}p{1cm}>{\centering\arraybackslash}p{1cm} >{\centering\arraybackslash}p{1cm} >{\centering\arraybackslash}p{1cm} >{\centering\arraybackslash}p{1cm} >{\centering\arraybackslash}p{1cm} >{\centering\arraybackslash}p{1cm} >{\centering\arraybackslash}p{1cm} >{\centering\arraybackslash}p{1cm} >{\centering\arraybackslash}p{1cm} >{\centering\arraybackslash}p{1.5cm}}
 
		\toprule
               &{\hspace{1em}} Method & \makecell[c]{\rotatebox{90}{Aircraft~\cite{maji2013fine}}} & \makecell[c]{\rotatebox{90}{Caltech101~\cite{fei2004learning}}} & \makecell[c]{\rotatebox{90}{CIFAR100~\cite{krizhevsky2009learning}}} & \makecell[c]{\rotatebox{90}{DTD~\cite{cimpoi2014describing}}} & \makecell[c]{\rotatebox{90}{EuroSAT~\cite{helber2019eurosat}}} & \makecell[c]{\rotatebox{90}{Flowers~\cite{nilsback2008automated}}} & \makecell[c]{\rotatebox{90}{Food~\cite{bossard2014food}}} & \makecell[c]{\rotatebox{90}{MNIST~\cite{deng2012mnist}}} & \makecell[c]{\rotatebox{90}{OxfordPet~\cite{parkhi2012cats}}} & \makecell[c]{\rotatebox{90}{Cars~\cite{krause20133d}}} & \makecell[c]{\rotatebox{90}{SUN397~\cite{xiao2010sun}}} & \makecell[c]{\textit{Average}} \\
  
		\midrule
		\multirow{3}{*}{\rotatebox{90}{CLIP}}&{\hspace{1em}}Zero-shot & 24.3 & 88.4 & 68.2 & 44.6 & 54.9 & 71.0 & 88.5 & 59.4 & 89.0 & 64.7 & 65.2 & 65.3  \\
            &{\hspace{1em}}5-shot Full Fine-tune & 30.6 & 93.5 & 76.8 & 65.1 & 91.7 & 92.9 & 83.3 & 96.6 & 84.9 & 65.4 & 71.3 & 77.5 \\
            &{\hspace{1em}}5-shot Fine-tune Adapter & 29.7 & 90.0 & 75.3 & 63.9 & 81.1 & 94.2 & 87.8 & 90.4 & 89.0 & 68.2 & 72.5 & 76.6 \\
            
            \midrule\midrule
            \multirow{7}{*}{\rotatebox{90}{\textbf{Transfer}}}&{\hspace{1em}}Continual-FT & & 72.8  & 53.0 &36.4  &35.4  & 43.3 & 68.4 & 47.4 & 72.6 & 30.0 &52.7  & 51.2   \\
            &{\hspace{1em}}LwF \cite{li2017learning} &  & 72.1 & 49.2 & 35.9 & 44.5 & 41.1 & 66.6 & 50.5 & 69.0 & 19.0 & 51.7 & 50.0   \\
            %{\hspace{1em}}{\hspace{1em}}iCaRL \cite{rebuffi2017icarl} & \textbf{\textcolor{red}{running}} &  &  &  &  &  &  & &  &  &  &  \\
            &{\hspace{1em}}LwF-VR \cite{ding2022don} &  & 82.2 & 62.5 & 40.1 & 40.1 & 56.3 & 80.0 & 60.9 & 77.6 & 40.5 & 60.8 & 60.1   \\
            &{\hspace{1em}}WiSE-FT \cite{wortsman2022robust} &  & 77.6 & 60.0 & 41.3 & 39.4 & 53.0 & 76.6 & 58.1 & 75.5 & 37.3 & 58.2 & 57.7   \\
            &{\hspace{1em}}ZSCL \cite{zheng2023preventing} &  & \underline{84.0} & \underline{68.1} & \textbf{44.8} & \underline{46.8} & \underline{63.6} & \underline{84.9} & \underline{61.4} & \underline{81.4} & \underline{55.5} & \underline{62.2} & \underline{65.3} \\
            \rowcolor{Thistle!20}
            &{\hspace{1em}}Ours &  & \textbf{87.9} & \textbf{68.2} & \underline{44.1} & \textbf{48.1} & \textbf{64.7} & \textbf{88.8} & \textbf{69.0} & \textbf{89.1} & \textbf{64.5} & \textbf{65.1} & \textbf{68.9\textcolor{Maroon}{(+3.6)}} \\
            % ----------------------------------------------------------------------
            \midrule
            \multirow{7}{*}{\rotatebox{90}{\textbf{Average}}}&{\hspace{1em}}Continual-FT\hphantom{Space} \hphantom{Space} & 28.1 & 86.4 &59.1 &52.8 & 55.8 &62.0 &70.2 & 64.7 & 75.5 & 35.0 &54.0 & 58.5 \\
            &{\hspace{1em}}LwF \cite{li2017learning} &23.5 &77.4  &43.5 &41.7 &43.5  &52.2 &54.6 & 63.4 & 68.0& 21.3& 52.6&49.2\\
            %{\hspace{1em}}{\hspace{1em}}iCaRL \cite{rebuffi2017icarl} & \textbf{\textcolor{red}{running}} &  &  &  &  &  &  & &  &  &  &  \\
            &{\hspace{1em}}LwF-VR \cite{ding2022don} &24.9 & \underline{89.1} &64.2 &53.4 &54.3  &70.8 &79.2 &66.5  &79.2 & 44.1& 61.6&62.5\\
            &{\hspace{1em}}WiSE-FT \cite{wortsman2022robust} & \textbf{32.0} & 87.7 & 61.0 &\underline{55.8} & \underline{68.1} & 69.3&76.8 &\underline{71.5}  &77.6 &42.0 &59.3&63.7  \\
            &{\hspace{1em}}ZSCL \cite{zheng2023preventing} & 28.2& 88.6 &\underline{66.5} & 53.5&56.3  &\underline{73.4} &\underline{83.1} & 56.4 & \underline{82.4} & \underline{57.5}&\underline{62.9}&\underline{64.4} \\
            \rowcolor{Thistle!20}
            &{\hspace{1em}}Ours & \underline{30.0} & \textbf{89.6} & \textbf{73.9}& \textbf{58.7}& \textbf{69.3} &\textbf{79.3} & \textbf{88.1}& \textbf{76.5} & \textbf{89.1}& \textbf{65.3}&\textbf{65.8} & \textbf{71.4\textcolor{Maroon}{(+7.0)}} \\
            
            \midrule
            \multirow{7}{*}{\rotatebox{90}{\textbf{Last}}}&{\hspace{1em}}Continual-FT & 27.8 & 86.9 & 60.1 & 58.4 & 56.6 & 75.7 & 73.8 & \underline{93.1} & 82.5 & 57.0 & 66.8 & 67.1 \\
            &{\hspace{1em}}LwF \cite{li2017learning} & 22.1 & 58.2 & 17.9 & 32.1 & 28.1 & 66.7 & 46.0 & 84.3 & 64.1 & 31.5 & 60.1 & 46.5 \\
            %{\hspace{1em}}{\hspace{1em}}iCaRL \cite{rebuffi2017icarl} & \textbf{\textcolor{red}{running}} &  &  &  &  &  &  &  &  &  &  &  \\
            &{\hspace{1em}}LwF-VR \cite{ding2022don} & 22.9 & \textbf{89.8} & 59.3 & 57.1 & 57.6 & 79.2 & 78.3 & 77.7 & 83.6 & 60.1 & 69.8 & 66.9  \\
            &{\hspace{1em}}WiSE-FT \cite{wortsman2022robust} & \textbf{30.8} & 88.9 & 59.6 & \underline{60.3} & \underline{80.9} & 81.7 & 77.1 & \textbf{94.9} & 83.2 & 62.8 & 70.0 & \underline{71.9}  \\
            &{\hspace{1em}}ZSCL \cite{zheng2023preventing} &26.8 & 88.5 & \underline{63.7} & 55.7 & 60.2 & \underline{82.1} & \underline{82.6} & 58.6 & \underline{85.9} & \underline{66.7} & \underline{70.4} & 67.4 \\

            \rowcolor{Thistle!20}
            &{\hspace{1em}}Ours & \underline{30.1} & \underline{89.3} &\textbf{74.9}  & \textbf{64.0} & \textbf{82.3} &  \textbf{89.4}&\textbf{87.1} & 89.0 & \textbf{89.1} &\textbf{69.5} &\textbf{72.5}  & \textbf{76.1\textcolor{Maroon}{(+4.2)}} \\ 
            
		\bottomrule
	\end{tabular}}
 \vspace{-4pt}
	\caption{Comparison with state-of-the-art methods on few-shot MTIL benchmark in terms of ``Transfer”, ``Average”, and ``Last” scores (\%). Ours converges in 500 iterations on few-shot. We label the best and second methods with \textbf{bold} and \underline{underline} styles. The top block indicates the upper-bound solutions to adapt the CLIP on each task.}
	\label{tab:comparefs}
  \vspace{-11pt}
\end{table*}

\vspace{1ex}
\noindent\textbf{Training MoE-Adapters.}
We train the MoE-Adapters through  simple back-propagation, orchestrated by an incremental activate-freeze strategy. The objective is to augment experts with intra-task knowledge and inter-task collaboration. 
Specifically, after training on an older task, we count the distribution of its router's outputs. 
The $Top$-$k$ most activated experts are then kept frozen during subsequent task training to preserve task-specific knowledge. 
In this manner,  when faced with a new task, the respective router is able to access frozen experts for leveraging shareable knowledge from historical tasks, and optimize unfrozen experts to acquire specific information for the new task.
As illustrated in Figure~\ref{fig:different frozen}, during training the router can activate (a) only the untapped experts, (b) both untapped and previously learned experts, and (c) only the learned experts from previous tasks. As a result, this strategy allows experts to consolidate their knowledge collaboratively, resembling the human brain's mechanism of reinforcing and connecting new information with existing memories.


%
%%%%%%%%%%%%%%%%%%%%%%%%%%%%%%%%%%%%%%%%%%%%%%%%%%%%%%%



%---------------------------------------------------------------------------------------------






\subsection{Distribution Discriminative Auto-Selector}
\label{DDAS}
The task-specific nature of the routers in our MoE-Adapters necessitates manual task identity to activate the appropriate router. Such manual intervention is not aligned with the automated and practical nature of Task Incremental Learning (TIL) and Class Incremental Learning (CIL), and restricts the inherent zero-shot generalization capability of the CLIP model. To address this limitation, we develop the Distribution Discriminative Auto-Selector (DDAS), which automatically selects the proper router with the task context inferred by analyzing the variation in the distribution of input.

DDAS extends the incremental MoE framework by introducing a series of task-specific autoencoders, $\{\mathcal{F}_A^t\}_{t=1}^T$, which are trained to independently capture the distribution characteristics for the tasks $\{{\mathcal{T}^t}\}_{t=1}^{T}$. The loss function employed for training the autoencoders is the Mean Squared Error (MSE), defined as:
\begin{equation}
     %d^t=MSE( \textit{\textbf{f}}^t_i, \textit{\textbf{f}}^t_o ),
     d^t= ||\textit{\textbf{f}}^t_i - \textit{\textbf{f}}^t_o||^2,
  \label{eq:chooseid1}
\end{equation}
where $\textit{\textbf{f}}^t_i$ is the intermediate feature extracted from the input image, and $\textit{\textbf{f}}^t_o=\mathcal{F}_A^t(\textit{\textbf{f}}^t_i)$ is the reconstructed feature representation by the autoencoder of task $t$. Since the autoencoder $\mathcal{F}_A^t$ is individually learned on the data of task $t$, the resulting reconstruction score $d^t$ reflects the likelihood that an input image pertains to the task, with a lower score suggesting a higher probability.

Moreover, to preserve CLIP's zero-shot transfer ability during continual learning, we include an additional autoencoder, $\mathcal{F}^0_A$, trained on a reference dataset to identify out-of-distribution data. Upon completion of the learning process, DDAS computes a set of distribution scores $\{d^t\}_{t=1}^T$ for each input image. Should all scores surpass a specific threshold, $Thres$, the system classifies the input as ``unseen data" and redirects it to the frozen CLIP for zero-shot transfer. Conversely, inputs below this threshold are routed to the corresponding router with the lowest distribution score, ensuring efficient and accurate task identification.

%  
